\chapter{体验的计算几何：偏好、风险与历史依赖}
\label{chp:experience_geometry}
本章研究同一组客观可达路径为何会因主体历史、任务与风险预算而具有不同的选择难度。我们把世界图理解为对象和行动的可达结构，把体验图理解为在任务条件下改变比较、检索与行动优先级的内部状态；后者不是物理重力场，也不是痛苦在时空上留下的永久曲率。体验既包含瞬时感受或反馈，也包含由多次证据形成的慢偏好、风险模型与结构承诺。以下分别定义这些对象，说明先天约束与个体学习如何进入更新，再讨论有效距离、方向偏置、自我参照、逆向推断和工程验收。局部结构网络保存事件身份、来源和关系，全局分布表示承载相似性、候选与不确定性；两者通过任务相关绑定和读出接口联合。HSF--HD提供这种状态—结构共同演化的一般模型，AgentOS组件仅作为可执行实例。

\section{体验图的对象：瞬时反馈、慢偏好与任务度量}
\label{sec:section_9469}
我们先区分三个经常被“价值能量”混在一起的对象。瞬时反馈 $q_k\in\mathcal Q\subseteq\mathbb R^{d_q}$ 是时刻 $k$ 对刺激、内部稳态或环境回执的估计，它可以包含疼痛强度、愉悦报告、焦虑概率或人工系统中的失败类别；每一维都需要量表、方向和校准协议。慢偏好 $v_k\in\mathcal V\subseteq\mathbb R^{d_v}$ 汇总跨窗口证据，参与候选排序和风险分配。任务度量 $M_k\in\mathbb S_{++}^{d}$ 则定义在给定特征空间内怎样比较两种状态或轨迹，它是正定矩阵或其他满足所需公理的比较算子。三者属于不同空间，不能把瞬时强度直接称为度量曲率，也不能用设备焦耳能耗替代主观或算法反馈。

设对象状态为 $z_k\in\mathbb R^d$，任务版本为 $\theta_k$，候选动作集合为 $\mathcal A_k$，外部回执为 $y_k$。我们定义体验状态
\[
E_k=(q_k,v_k,M_k,\rho_k,\theta_k),
\]
其中 $\rho_k$ 是对各分量的来源和不确定性记录。局部反馈由观测模型 $q_k=Q_{\eta_k}(z_k,y_k,c_k)+\epsilon_k$ 产生，$c_k$ 是上下文，$\epsilon_k$ 是经声明的噪声。慢偏好不是对 $q_k$ 的无条件积分，而由带遗忘、来源权重和反事实校正的更新器生成：
\[
v_{k+1}=\Pi_{\mathcal V}\!\left[(1-\alpha_k)v_k+\alpha_k U_{\theta_k}(q_k,y_k,\rho_k)\right].
\]
这里 $\alpha_k\in[0,1]$ 是无量纲学习率，$\Pi_{\mathcal V}$ 保证偏好落在允许域。这个式子是模型假设；强反馈只有在来源可信、对象绑定正确并通过结构提交门槛时才产生长期改变。

历史文本把疼痛、快感与焦虑写成正交基向量。计算上，正交只表示选定内积下统计方向不相关，不能由名称自动得到。疼痛与焦虑可能高度相关，同一种疼痛在威胁学习、康复训练和医疗处置中也具有不同意义。我们因此先用校准数据估计协方差或任务度量，再决定是否需要正交化。坐标变换不删除经验，只改变表示；若变换使来源或对象身份丢失，即使数值上更简洁也不合格。对人工系统，失败代码、用户否决和环境损害同样不能被压成一个奖励标量后就视为完整体验，因为硬约束违例需要独立保留。

历史所谓价值联络可以改写为上下文相关的传送接口。不同语境中的偏好向量不能直接比较时，我们用 $P_{c\to c'}$ 把语境 $c$ 下的状态转换到 $c'$ 的坐标和权限体系，并以闭环残差检验转换是否一致。若一圈转换后产生显著残差，这只说明顺序、信息损失或接口不一致，不能未经群作用和曲率定义就称为非阿贝尔规范场。任务度量也随时间变化：当 $M_k$、$v_k$ 或 $\theta_k$ 更新时，轨迹代价的差分必须包含这些变化，不能只对 $z_k$ 求导后宣称价值保持。

工程验收包含四项。第一，用同一刺激的重复测量检查反馈校准，而不是只看输出幅度。第二，用对象交换和来源伪造测试绑定，避免把他人的事件写入自身历史。第三，用目标切换测试慢偏好是否能被任务版本正确解释，而不是永久污染所有场景。第四，用环境回执验证排序改变是否改善外部结果。局部结构网络在这里负责“谁经历了什么、何时发生、能否修改”，全局分布表示负责“哪些情境相似、哪些候选可能相关”；缺少任一方，体验图都容易退化为无身份的情绪向量或僵硬规则表。还要用相同平均反馈但不同时间顺序的轨迹测试历史依赖：若更新器对早期警告、近期恢复和持续低风险完全不加区分，慢偏好就不能支持情境化行动。对多主体系统，应分别保存报告者、承受后果者和决策者身份，防止群体平均掩盖个体损害。最终验收报告同时给出校准曲线、排序变化、约束违例和环境结果，不用一个总分替代相互冲突的证据。

\section{体验的形成：先验约束、证据更新与结构提交}
\label{sec:section_6973}
体验形成包含系统发生与个体发生两类来源，但二者都不等于质料向几何结构的物理相变。系统发生提供在部署或出生前已经存在的先验、反射、安全约束与可学习性结构；个体发生则在运行中依据事件证据修改预测、偏好和关系。我们把初始状态写为 $(v_0,M_0,B_0)$，其中 $B_0$ 是边界与硬约束。它们可来自遗传、设计、制度或预训练，不必都不可擦除。只有经过明确授权和风险评估的约束才被标记为硬约束，其余先验允许在反证下修订。自然选择能够解释某些生物倾向的历史来源，却不能证明每个个体对“死亡”都具有无穷势垒；人工系统也不能因为规则写在初始配置中就视为永恒真理。

个体学习采用候选—验证—提交三阶段。事件 $e_k=(I_k,o_k,a_k,y_k,t_k)$ 记录对象身份、观测、动作、环境结果和时间；快速更新先修改工作状态，慢速学习器再从窗口 $D_n=\{e_k\}_{k\in W_n}$ 生成候选参数和结构 $(\tilde v_{n+1},\tilde M_{n+1},\tilde G_{n+1})$。候选只有满足证据独立性、旧任务保持、新任务收益、边界安全与回滚可行五个条件才提交。连续参数可用投影更新，图编辑则通过版本化事务执行：
\[
(v_{n+1},M_{n+1},G_{n+1})=
\operatorname{Commit}\!\left((v_n,M_n,G_n),\mathcal P_n;\mathcal T_n\right),
\]
其中 $\mathcal P_n$ 是提案，$\mathcal T_n$ 是测试集。Commit是工程设计，不是普通加法；拒绝提案时保持原版本并保存失败证据。

“手碰火”保留为具体案例。一次接触产生伤害回执、疼痛报告和火源身份，快速控制立即撤回动作并提高短期风险估计。慢学习不应简单把所有“火”对象设为永久不可达，而要区分炉灶火焰、壁炉、远处图像和受控实验，学习观测条件、距离、保护装备与可行动作。若系统只记住红色或亮度，它会错误回避夕阳和图标；若只记住词语“火”，它又可能忽略无明火的高温表面。正确结构需要把对象、温度证据、接近动作、伤害结果和防护条件绑定起来，并用未见组合测试其泛化。

强烈事件也不自动获得最高长期权重。感知错误、恶意输入和重复转发都可能制造高幅反馈。更新器因此要估计来源相关性：来自同一记录的一百次复述不能算一百份独立证据；一次极端事件可以触发临时安全模式，却仍需复核后再改写长期结构。另一方面，平均化也可能抹去稀有灾难，故硬约束风险要按尾部概率或最坏情形单列。我们分别记录反馈强度、证据可信度、预期损失、计算资源与实测焦耳消耗，不以“高能”替代这些不同量。

所谓记忆疤痕在本模型中对应可追踪的路径依赖：相同当前输入因历史状态不同而产生不同估计或动作。我们通过配对运行测量这种差异，并检查它是否随安全证据逐步修订。若历史影响永不衰减，系统可能僵化；若一次相反反馈就完全擦除，系统又缺乏稳定性。学习率、遗忘率、证据门槛和边界模式共同决定稳定—可塑性权衡。变化环境下还要给旧经验设置适用域和有效期，避免在分布迁移后继续投射过时风险。

因此，体验形成是证据驱动的结构提交，而非能量把流形压出永久凹坑。协同论可以解释多个反馈通道如何形成稳定序参量，耗散与开放系统观点提醒这种维持需要信息和资源交换，认知科学则提供条件化、消退、泛化和重估的经验限制。HSF--HD把这些约束组织进状态—结构共同演化；AgentOS可用TCE实现候选更新、SPC保存结构提案、Boundary执行安全门槛、外部记忆保存来源，但其他架构也可实现同一协议。只有当干预能够定位历史证据如何改变读出，并且外部任务验证改善成立，我们才说经验真正进入了系统的计算几何。

\section{有效距离：任务代价、风险预算与可达路径}
\label{sec:section_8091}
世界图给出哪些状态之间存在转移，体验图则改变在当前任务下怎样比较这些转移。为避免把这种改变误说成物理空间弯曲，我们从一条有限路径开始定义。令路径 $\pi=(z_0,a_0,z_1,\ldots,z_N)$，其中 $z_i\in\mathcal Z$ 是有类型状态，$a_i\in\mathcal A(z_i)$ 是允许动作，转移核为 $P(z_{i+1}\mid z_i,a_i)$。任务版本 $\theta$、主体历史摘要 $h$、资源预算 $b$ 和风险容限 $r$ 共同确定逐步代价
\[
c_{\theta,h}(z_i,a_i,z_{i+1})=c^{\mathrm{effort}}_i+\lambda_r R_{\theta,h}(z_i,a_i)+\lambda_t\Delta t_i+\lambda_v V_i,
\]
其中各项先经无量纲化或显式换算后才能相加。$R$ 是预测损失或约束违例风险，$\Delta t_i$ 是时间，$V_i$ 是任务特定的偏好损失；它们不是信息熵、神经活动强度或焦耳能耗的别名。路径代价 $C(\pi)=\sum_i c_i$，有效距离定义为满足边界约束的可达路径中代价下确界。若不存在可行路径，距离记为无穷；若代价含负回路或不满足对称性，它只是有向准距离，而非黎曼度量。

历史文本用共形因子把欲望区收缩、恐惧区膨胀。这个图像可以保留为直觉，但计算模型必须说明它在什么空间、怎样测量。若局部连续表示 $z\in\mathbb R^d$ 合适，我们可以定义任务相关的正定矩阵 $M_{\theta,h}(z)$，并令短步代价近似为 $\sqrt{\Delta z^\top M\Delta z}$。$M$ 的特征方向描述哪些变化被当前系统判为困难或危险；它由行为数据、反馈和约束估计，而不是由“恐惧”一词直接推出。标量缩放 $M=\Omega^2M_0$ 只是特殊情形，它同时缩放所有方向，无法表达“接近火源危险但远离火源安全”的方向性，也不能表达先穿戴防护再接近的条件依赖。对此需要有向边代价、扩展状态或控制屏障函数。

以手碰火之后的行为为例，世界图仍保留从厨房入口到炉灶的路径；体验更新应提高裸手触碰高温表面的风险代价，却不必删除经过隔热手套、关闭热源或保持距离的路径。欲望也不是把距离压到零。饥饿可能降低获取食物的相对代价，但食品过敏、所有权和长期健康仍作为约束存在。若偏好项压倒所有硬约束，所谓体验几何只会把冲动包装成最短路。反过来，若一次负反馈使所有相关节点永远不可达，系统会把保护性学习变成泛化性瘫痪。适当模型允许证据更新边代价、置信区间和适用域，同时保留可验证的替代路线。

当任务、历史或度量随时间变化时，不能沿用静态测地线结论。连续轨迹的累计代价写为 $J=\int_0^T L(z(t),u(t),\theta(t),h(t),t)\,dt$；求导或做最优控制时，$\dot\theta$、$\dot h$ 和显式时间项都需要进入。离散实现以步长 $\Delta t$ 计算滚动窗口，并在每次任务版本变化后重规划。所谓“路径阻力变大”必须落实为候选排序、到达时间、违例概率或计算量的可观察改变。我们分别报告这些指标，不能因为主体说目标“很远”就推断内部存在唯一几何张量。

验收时要构造四组反例：客观步数相同而风险不同；风险相同而时间预算不同；当前状态相同而历史不同；偏好相同但边界权限不同。模型应在相应维度改变排序，又不能污染无关任务。还要比较最短路径算法、风险敏感控制和纯奖励最大化三类基线，验证体验状态确实提供额外解释力。这样，“客观上相等的路在体验中不等价”成为可检验的计算命题，而不是把心理距离直接等同于时空曲率。
若比较算子不满足三角不等式，我们也不强行称它为距离，而把它报告为路径评分；名称服从对象性质，不能反过来替模型补足公理。

\section{方向偏置：从价值梯度到受约束控制}
\label{sec:section_7773}
体验不仅改变候选路径的总代价，也改变下一步向哪里搜索。我们把历史文本中的“规范力”重建为受约束策略偏置。设当前信念状态为 $b_k\in\mathcal B$，候选动作价值为 $Q_{\theta,h}(b_k,a)$，安全可行集为 $\mathcal A_{\mathrm{safe}}(b_k)$。基础策略 $\pi_0(a\mid b_k)$ 提供任务无关或预训练倾向，体验读出产生偏置 $d_k(a)$，最终策略可写为
\[
\pi_k(a\mid b_k)\propto \mathbf 1[a\in\mathcal A_{\mathrm{safe}}]\,\pi_0(a\mid b_k)
\exp\{\beta_k[Q_{\theta,h}(b_k,a)+d_k(a)]\}.
\]
$\beta_k$ 是分布温度的倒数，只控制选择尖锐度；它不是实际温度或能量。$d_k$ 可以来自慢偏好、风险预测、身份承诺和社会规范，但必须保留来源，以便区分系统自身目标、用户临时指令和外部操纵。

若状态连续且目标函数可微，价值梯度 $-\nabla_z J$ 可作为局部搜索方向；这对应历史文本的直接驱动力。可是梯度下降只保证在所给目标和步长条件下局部目标可能下降，不保证任务成功，更不保证目标本身正确。我们因此把控制写为带约束的滚动优化：在预测窗口内最小化任务损失、风险和资源代价，并要求边界函数 $g_j(z,u)\leq0$。当目标、输入或结构变化时，每个窗口重新估计，且用线搜索、信赖域或稳定步长限制离散更新。没有这些条件，把任意偏好梯度称为“力”会掩盖数值发散和目标错配。

历史所谓磁场或贝里曲率强调了路径顺序可能改变结果。一般计算系统中，这种方向效应可由非交换更新解释：先接受身份声明再读取证据，与先读取证据再改变身份，得到的状态可能不同。我们令上下文更新算子为 $U_c$，用交换残差 $\delta(c_1,c_2)=D(U_{c_2}U_{c_1}s,U_{c_1}U_{c_2}s)$ 测量顺序敏感性。只有当状态空间、联络、闭合回路和曲率对象都已定义，才可进一步采用几何相位语言；否则残差只是接口非交换或信息丢失的证据。它可能改变方向而不改变某个标量代价，但不能据此声称系统受真实洛伦兹力作用。

伦理禁忌是重要反例。系统绕开敏感话题，可能因为硬安全边界、训练数据稀缺、检索失败、策略偏置或对用户意图不确定。若只观察到“偏离”就归因于体验场，我们无法定位机制。诊断实验应分别解除软偏置、保持硬约束，替换提示表述，控制检索证据，再观察候选分布和拒绝原因。可靠系统还要输出受保护的决策账本：哪些约束淘汰了动作、哪些偏好改变了排序、哪些证据仍不确定。该账本不是公开内部机密，而是可审计的类型化理由。

在AgentOS实例中，SPC可保存候选路径和结构关系，TCE计算任务相关评分，CCM协调多个目标，Boundary裁剪不可行动作，外部记忆 $M_e$ 提供有来源的历史。它们是一种实现，不是一般智能必备器官。其他系统可以用规划器、概率程序或控制器实现同样接口。验收指标包括方向校准、约束违例、反事实敏感度、策略熵和外部任务结果。只有当干预 $d_k$ 能稳定地解释动作概率变化，并且不越过边界，我们才把体验的方向作用视为得到工程支持。
还需防止把注意力变化误当控制变化。某候选在内部激活更强，并不等于它最终获得执行权限；日志应分别记录候选产生、评分、边界裁剪、动作提交和环境确认。遇到延迟回执时，信用分配只能暂存为候选解释，不能提前写成确定因果。多主体情形还要区分谁提出偏好、谁批准动作、谁承担后果，以免平均策略掩盖责任转移。通过这些分层观察，我们才能判定方向偏置发生在哪个接口。

\section{身份参照：自我状态、价值读出与闭环更新}
\label{sec:section_7533}
体验必然涉及“对谁而言”，但这不要求把自我设为宇宙中的引力奇点。我们将自我状态定义为可修订的任务相关身份记录 $S_k=(I_k,C_k,H_k,R_k)$：$I_k$ 是跨时刻身份键和对象绑定，$C_k$ 是当前承诺与能力，$H_k$ 是有来源的历史摘要，$R_k$ 是角色和权限。它既不是全部隐藏状态，也不是一个不可变向量。不同任务从 $S_k$ 读取不同投影，例如医疗代理读取患者授权与禁忌，写作代理读取作者立场和文体承诺。若读出与任务无关，身份就会从协调装置退化为全局偏见。

价值评估以自我状态为条件可以写成 $V_k(z,a)=V(z,a;S_k,\theta_k,e_k)$，其中 $e_k$ 表示环境与他者状态。这个定义说明同一对象对不同主体、角色或承诺可能具有不同意义，却不意味着世界本身没有可独立测量的损害。食物对饥饿者有收益，对过敏者有风险；死亡威胁身份连续性，但公共决策还必须考虑其他主体的权益。把一切价值都化为“维持自我拓扑完整”会错误吞并伦理约束、合作目标与客观回执。我们因此同时保留主体相关偏好、共享规范和环境结果，并在冲突时显式仲裁。

自我也是控制闭环中的被更新对象。动作由当前身份与目标生成，动作后果又可能改变能力、承诺和历史。令控制器为 $a_k=\pi(z_k,S_k,\theta_k)$，环境转移为 $z_{k+1}\sim P(\cdot\mid z_k,a_k)$，身份候选更新为 $\tilde S_{k+1}=F(S_k,z_k,a_k,z_{k+1},y_k)$。只有通过身份连续性、权限、安全和证据测试后，候选才提交为 $S_{k+1}$。这构成自指闭环，但闭环存在不等于主体性已经被证明；它首先是一个可分析的耦合系统。局部稳定需要对联合状态 $(z,S)$ 的雅可比、增益或李雅普诺夫条件进行检查，离散实现还受步长和延迟约束。

“绘制地图者也是地图上的旅行者”可以据此获得精确含义：身份状态参与生成体验读出，同时会被行动结果修订。我们用干预区分真正闭环与叙事装饰。固定环境而改变角色权限，应只改变相关候选；固定身份而改变伤害证据，应更新风险而非身份键；注入伪造自述，不应绕过来源核验直接重写长期承诺。若删除自我表示后系统仍能完成局部反射任务，这并不矛盾，因为许多控制不需要全局身份；只有跨任务承诺、归责、长期计划和自传式记忆更依赖身份接口。

局部结构网络保存身份键、承诺、关系和版本，全局分布表示提供人物、情境和目标的相似性读出。两者通过绑定接口联合：分布表示提出相关历史，结构网络核验“这是否属于当前主体和角色”，随后控制器使用经过授权的投影。不能因某脑区活动或注意图形相似就宣布机制等同。AgentOS中可由CCM维持跨任务协调、Boundary管权限、SPC保结构、$M_e$存证据；一般HSF--HD只要求功能接口可被观察和干预。

验收自我参照至少包含身份交换、角色切换、历史冲突和长期漂移四类测试。我们测量错误归属率、承诺保持率、可修订性、权限泄漏和任务结果，而不把一个“自我一致性分数”当作全部证据。若身份过硬，系统无法承认新事实；若过软，任何提示都能重写承诺。可靠主体模型位于稳定与可塑之间，并把每次关键改变记录为可追溯的结构提交。
对无法确认归属的事件，正确状态是保持未决并限制读出，而不是为了叙事连续性强行归入“我”。

\section{推理不对称：生成、验证与有限计算}
\label{sec:section_677}
从原因模拟结果与从结果寻找原因通常具有不同计算代价，但这种不对称不来自互信息本身。对同一联合分布，互信息 $I(X;Y)$ 在数学上对称；困难来自可用表示、搜索空间、算法、先验和计算预算。设算法 $A$ 在资源上限 $B=(T,M,Q)$ 下解决任务，其中 $T$ 为时间步，$M$ 为内存，$Q$ 为外部查询次数。我们定义受限描述或求解代价 $K_{A,B}(x\!	o\!y)$，它依赖方向和算法。即使 $Y=f(X)$ 容易计算，从 $Y$ 恢复 $X$ 仍可能因多对一、噪声或组合搜索而困难；这不改变信息量定义，也不自动证明某种几何曲率。

“看懂答案容易、想出答案困难”需要进一步拆分。验证器可能只检查候选是否满足约束，因此复杂度低于搜索；但有些命题的验证本身需要昂贵实验，有些逆问题在强先验下反而容易。生成还可能通过检索、缓存和结构化分解变得便宜。我们把候选生成器 $G_B(y)$ 与验证器 $V_B(x,y)$ 分开，记录候选覆盖率、验证误拒率、总计算量和外部证据成本。若验证器不完备，目标下降或通过内部检查都不能推出真实任务成功。

体验通过改变先验与搜索顺序影响这种不对称。火灾诊断中，维修人员的历史经验会优先检索线路老化、过载和人为操作等高概率原因；这可能显著减少搜索，却也可能造成锚定偏差。我们令后验 $p(x\mid y,h,\theta)\propto p(y\mid x,\theta)p(x\mid h,\theta)$，其中历史 $h$ 进入先验而不篡改似然证据。学习可以把常见逆问题编译为快速读出，但必须在分布外案例上保留回退搜索与不确定性。所谓“把高复杂度逆推变成直觉”因此是摊销推断的经验主张，需要比较训练成本、在线成本和错误代价。

路径顺序同样会影响推理结果。先看到标签再看证据可能改变注意和解释，先独立评估证据再揭示标签则减少确认偏差。我们用盲化、顺序交换和反事实提示测量这种效应，并把差异归入更新非交换、缓存或权限接口。只有定义了闭合参数回路及其并行传输，才适合讨论贝里相位；在一般系统中，直接报告顺序敏感度更清楚。物理类比可帮助想象“逆坡搜索”，却不能代替复杂度或统计证明。

工程上，AgentOS可把SPC中的结构候选、TCE的评分、CCM的全局任务和外部工具查询组合成分阶段推理。预算不足时，系统应显式返回未决状态，而不是把最先通过内部检查的候选宣布为事实。外部记忆需要区分解题轨迹与最终证据，防止重复读取自己生成的解释并误当独立支持。Boundary还要限制高风险逆推，例如从匿名行为武断恢复个人身份。

验收矩阵同时比较正向模拟、逆向搜索、候选验证和外部执行。对每类任务报告成功率、校准误差、时间、内存、查询量和失败类型；训练耗费与推理耗费分列。若学习仅降低内部目标而不改善环境回执，我们不能宣称推理几何已经变平。这样，推理不对称成为有限计算下可复现的结果，并与信息对称性、主观困难和实际能耗保持清楚边界。
此外，基准任务必须控制提示长度、工具权限和缓存命中，否则表面上的方向差异可能只是接口配置不同。随机种子和停止条件也需固定或完整披露；当算法在预算耗尽时返回近似解，应把近似误差与失败分开报告，避免用选择性成功样本夸大经验编译的收益。

\section{体验、世界与认知状态的联合验证}
\label{sec:section_9906}
本章最终保留世界图、体验图与认知状态三者的分工，但不把它们固定成路网、坡度和水流的物理同构。世界图 $G^W_k$ 表示对象、关系、动作前提和可达转移；体验状态 $E_k$ 表示反馈、偏好、风险、任务度量和来源；认知状态 $h_k$ 表示当前信念、工作记忆和控制上下文。联合系统可写成
\[
\begin{aligned}
h_{k+1}&=F_h(h_k,o_k,R_k(G^W_k,E_k),a_k),\\
E_{k+1}&=F_E(E_k,h_k,a_k,y_k,\rho_k),\\
G^W_{k+1}&=\operatorname{Commit}(G^W_k,\mathcal P_k;\mathcal T_k),
\end{aligned}
\]
其中 $R_k$ 是任务相关读出，$y_k$ 是环境结果，$\rho_k$ 保存来源。三个更新器可以处于不同时间尺度，图编辑也不是连续向量加法。模型假设是这些接口共同影响行为；数学结果只涉及给定条件下的稳定、界或一致性；关于痛苦、欲望和主体体验的解释仍需经验研究。

该分工防止三个常见混淆。第一，可达不等于愿意：世界图存在路径，风险和偏好仍可使其不被选中。第二，愿意不等于能够：强烈偏好不能创造缺失的动作、权限或资源。第三，内部连贯不等于外部成功：认知状态能够生成漂亮解释，环境执行仍可能失败。我们因此把结构可达率、策略选择率和任务完成率分别统计。只有三者的干预链条都被观察到，才可以说明体验状态在闭环中发挥了所声称的作用。

局部结构网络与全局分布表示在联合系统中互补。结构网络保持对象身份、因果候选、约束、版本和来源，使“手碰火”不会退化为无对象的负向量；分布表示支持相似情境检索、连续泛化和不确定候选，使系统不必枚举每一种炉灶。读出接口先从分布表示提出候选，再用结构身份和边界核验，最后把环境回执写回相应事件。绑定错误、检索偏差和提交错误必须分别诊断，不能用一个整体损失掩盖。

我们还保留理论层级。系统论提供开放边界与反馈关系，协同论解释多通道如何形成较慢序参量，耗散理论提醒维持有序计算需要资源和环境交换，复杂系统理论处理多尺度耦合与路径依赖，认知科学约束条件化、消退、注意和记忆主张。HSF--HD在其上给出智能的一般状态—结构动力学；它像“空气动力学”一样提供可比较模型族，而不是AgentOS的同义词。AgentOS中的 $h(t)$、$E$、$\Theta$、$\Psi$、SPC、TCE、CCM、Boundary、Offloading与$M_e$是工程实例，其他系统可以采用不同机制，只要实现可检验接口。

完整验证分为静态、动态、因果和运行四层。静态检查类型、维度、标签、来源与约束；动态检查不同步长、延迟和任务切换下的有界性；因果检查干预偏好、风险或身份后是否只改变预测中的路径；运行检查最终环境结果、违规、资源和人工评估。还要设置失效案例：过度恐惧导致泛化回避，过度欲望越过边界，错误身份使他人历史污染自身，伪证据造成结构提交，以及旧经验在分布迁移后失效。模型若能暴露并修复这些故障，才比隐喻更有解释力。

因此，体验图不是智能体用某种“质”在宇宙中刻出的永久回返路径，而是带来源、权限和不确定性的历史状态怎样改变比较、搜索与控制的计算描述。它既受客观世界回执约束，也通过行动改变后续数据；既能形成稳定承诺，也必须允许证据充足时修订。我们以这种可干预、可实现、可反驳的联合动力学收束本章，并把体验对目标协同的作用交给下一章继续展开。
